\documentclass[12pt]{article}


%%%%%%%%%%%%%%%%%%%%%%%%%%%%%%%%%%%%%%%%%%%%%%%%%%%%%%%%%%%%%%%%%%%%%%%%%%%%%%%%%%%%%%
%%%%%%%%%%%%%%%%%%% YOU PROBABLY DON"T NEED TO CHANGE ANYTHING HERE %%%%%%%%%%%%%%%%%%%
\usepackage[numbers,sort]{natbib}
\usepackage{amsfonts, amsmath} 
\usepackage{graphicx}
\usepackage{libertine} 
\usepackage{fullpage}
\usepackage{hyperref}
\hypersetup{
    colorlinks = true,
    linkcolor = blue,
    urlcolor  = blue,
    citecolor = blue,
    anchorcolor = blue}
\urlstyle{same}
%%%%%%%%%%%%%%%%%%%%%%%%%%%%%%%%%%%%%%%%%%%%%%%%%%%%%%%%%%%%%%%%%%%%%%%%%%%%%%%%%%%%%%
%%%%%%%%%%%%%%%%%%%%%%%%%%%%%%%%%%%%%%%%%%%%%%%%%%%%%%%%%%%%%%%%%%%%%%%%%%%%%%%%%%%%%%


%
% ===== DOCUMENT STARTS HERE =======
% you do not have to modify anything above
% if you do not wish to (but you indeed can0
%
\begin{document}

\title{Effects of Hyperparameters on Super Tiny Language Models (STLM)}
\date{\today}
\maketitle

{\bf Group Composition:}
\begin{enumerate}
\item Tan Yuvin, A0272629W, e1122032@u.nus.edu
\end{enumerate}

{\bf Github Repository:} \url{https://github.com/iNvy33/dsa4212-llm/tree/master}

%
%
\section{Introduction}


In this project, all runs were given 30min of wall-clock runtime on a 3060 TI. This is equivalent to roughly $3.0\times10^{16}$ FLOPs, which is order of magnitudes smaller than current small language models. For reference, the 124M version of one of the firsts open-weight model GPT-2 has an approximate FLOPs of $7\times10^{18}$ \url{https://github.com/karpathy/llm.c/discussions/481}. Our scope both in terms of data and model size must be small to cooperate with this limitation. 
\newline

The popular 'Tiny Stories' dataset on hugging face is perfect for tiny models, with its limited vocabulary size restricted to around 1500 words. Tiny Stories is made up of synthetically generated short stories from GPT-3.5 and GPT-4 meant to simulate language used in children's books. To supplement this artificial data with real life data, the "Children Books Text" dataset which contains children's book was also used when testing the models. 

\section{Baseline Model}

We first decide of the rough range parameter wise our model should sit in. The 3060 TI has a TFLOP/s of 32, translating to a maximum $5.76\times10^{16}$ FLOPs in 30min. How efficient the GPU is utilised during training depends on Model FLOPs Utilisation (MFU) which varies wildly, but using a value of 20\%, and the training compute = 6 x parameters x tokens, plus the Chinchilla ratio, we get within the ballpark of 9.8million parameters and $1.96\times10^8$ training tokens
\newline

We used BPE tokenization set to 4096 vocabulary size based on the total number of words in the dataset and the total training size to avoid repetition. A context size of 1024 was used, which covered over 99.9\% of the stories in the data set. The values which affect parameter size, which are the number of layers (8), token width (384), feed forward network width (768) were picked based on similarly sized models like nanoGPT. 
\newline

The AdamW optimizer with betas (0.9, 0.95) and weight decay 0.1 was used. In each configuration, a much shorter training sequence is ran, where in each iteration, the learning rate is increased. The plot of loss against iterations is used to estimate a suitable peak learning rate value for that particular model, shown in \autoref{fig:lr}. We detected the learning rate where the loss starts increasing, and take 1/3 of that value. The actual learning rate follows a cosine annealing with warm up shape, where it rises linearly to this peak in the first 2\% of the training, before following a cosine curve down to 10\% by the end.
\newline

\begin{figure}[h]
\centering
\includegraphics[width=0.8\textwidth]{Plots/baseline_lr.png}
\caption{Deciding optimal learning rate (red)}
\label{fig:lr}
\centering
\end{figure}

Other information about this model is that it uses, 4 headed attention, GeLU activation function, the more modern RMSNorm, applied pre-norm outside of the residual connection, and its batch size is 32768. 
\newline

In this report, several models will be compared, through their test Bits-per-byte(BPB) performance, which measures the average number of bits a model needs to predict each raw byte of text. The comparisons answers the following questions
\begin{enumerate}
\item With this training constraint, what model size gives the best result
\item Given a fix parameter number, which parameter allocations gives the best result
\item Can a simpler character based tokenization be used on this simple data, and how does this interact with training data size
\item Can ideas from mid-training in more modern LLMs be used to improve results
\item How does the model perform in a more general case
\end{enumerate}

\section{Results}

\subsection{Baseline Runs}

We first run the baseline model and observe its training loss over time in \autoref{fig:baseline} . While loss minimization is not fully monotone, due to us taking batches of data at a time, the overall trend is quite smooth. At this 8-layer model with 11.02M parameters, we can see that training loss just about starts converging towards the final minute of training. 
\newline

We ran the baseline model again with a different seed to get a rough estimate of stochastic variability within a model. \autoref{fig:baseline} shows the performance of each model on the Tiny Stories test set. The difference in BPB between the two runs of this model is 0.0009, which is by far the smallest among all the models. While 2 runs is not sufficient to get an accurate estimate of variance within models, we can be quite confident that any observed difference in model performances is due to the models themselves, and not any noise variance.

\begin{figure}[h]
\centering
\includegraphics[width=0.8\textwidth]{Plots/baseline.png}
\caption{Two identical runs with different seed}
\label{fig:baseline}
\centering
\end{figure}

\subsection{Parameter Size}

Next, we tested models of different parameter sizes. The simplest way to do so is to change the number of transformer block layers, where each block added an additional 1.18M parameters. The smaller models performed noticeably better at the start of the training, where at around 5 min of training, the smallest model containing 6.29M parameters had the smallest loss. Overtime, this smallest model quickly got overtaken by the other two. The largest model started overtaking the baseline model at around the 20min mark, though both stayed very close to each other thereafter. This could suggest that given a much larger training compute, larger models may still struggle to greatly outperform smaller models given this architecture. 

\begin{figure}[h]
\centering
\includegraphics[width=0.8\textwidth]{Plots/size.png}
\caption{4 Layers vs 8 Layers vs 12 Layers}
\label{fig:size}
\centering
\end{figure}

\subsection{Parameter Allocation}

Next, we took the total parameters in the baseline model, and assigned them different to the different parts of the transformer layer. We compared 2 additional models, assigning more parameters to either the token embedding width, or the feed forward network layer. \autoref{fig:allocation} shows that the rate of loss improvement is near identical, with the model that uses more parameters in the token embedding layer clearly performing worse than the other 2. The model with a larger FFN hidden layer performed marginally better by 0.002bpb, though this value is arguably within the noise interval, when referencing the 0.0009bpb difference in loss due to stochastic variability alone. 
\newline

An additional caveat is that the worser model had 11.25M parameters while the better one had 10.82M parameters. 

\begin{figure}[h]
\centering
\includegraphics[width=0.7\textwidth]{Plots/allocation.png}
\caption{Larger token representation width vs larger FFN hidden layer}
\label{fig:allocation}
\centering
\end{figure}

\subsection{Tokenization Choice}

Next, we compared this baseline model against a model that uses a much simpler character based tokenization instead. Keeping everything else the same bar the tokenization embeddings and context size, this character level tokenizer model contained 9.54M parameters. With smaller tokens, the model can only get through a much smaller portion of the total data in the 30 provided minutes. Since with this training time, the baseline BPE model does not use up the entire data (0.78 epochs), training data size is not a bottleneck. Given this, it is unsurprising that this model's performance is the worse yet.  
\newline

\begin{figure}[h]
\centering
\includegraphics[width=0.7\textwidth]{Plots/chr_BPE.png}
\caption{5\% vs full data training loss}
\label{fig:5pct}
\centering
\end{figure}

We now investigate the case when the training data could become such a bottleneck. We take 5\% of the total Tiny Stories dataset and discard the rest. At this data size, training for 30min uses about 15.5 epochs for the BPE model, where the model starts overfitting on the training data. On the other hand, the character level tokenization model uses up about 4.7 epochs, which is an optimal number of training data repeats to still maintain useful information. The jagged purple line shows that each time an epoch is used up, the model can predict the same data again more accurately, showing that it is remembering the training data.
\newline

We see that while the test BPB is still lower for the BPE model, there is a significant drop in performance, while the character level model maintained its BPB score. We see that the benefit of the BPE tokenizer is not just that the model can see more text, but also learns the most common letter combinations, whereas character level tokenization has to guess each byte within a BPE token. 

\begin{figure}[h]
\centering
\includegraphics[width=0.8\textwidth]{Plots/chr_inter.png}
\caption{Test results for character vs BPE tokenization}
\label{fig:inter}
\centering
\end{figure}

\subsection{Mid-Training}

Here, we borrow the idea "mid-training" from modern LLMs. The goal of mid-training, which comes after pre-training, is to include higher quality data more closely related to the task at hand. Since models remember latter training data better, assigning highly quality data towards the end of the training nudges the model to have a bias towards this better quality data.
\newline

From the Tiny Stories dataset page, stories generated from GPT-4 were deemed to be of higher quality. Thus, we re-arranged the training data, such that the first 80\% of the training was the same, but the latter 20\% was only on GPT-4 stories. We also tested on 3 different sets of validation data as shown in \autoref{fig:mid} . Based on the same validation set that was used throughout, this newer model actually performed quite a bit worse than the baseline. However, when testing on the other two sets of validation data, this mid-trained model performed very similarly across all sets, whereas, there is a huge difference in loss for the baseline model.
\newline

\begin{figure}[h]
\centering
\includegraphics[width=0.8\textwidth]{Plots/mid_training.png}
\caption{Different validation scores when partitioning the training data}
\label{fig:mid}
\centering
\end{figure}

From this test, we can conclude that the GPT-4 data is harder to predict, which could mean it is less simplistic and formulaic. Additionally, this mid-trained model is more general and can potentially be applied to a wider range of data.
\newline

To prove this second point, we compared both models on a brand new dataset, the Children's Book Test. The idea is that since the GPT generated stories were attempting to recreate the language used in real children's books, this dataset is representative of a more general test. Results show that this mid-trained model consistently outperformed the baseline model. A possible caveat is that these children books contain words unseen from the GPT-3.5 generated stories, but present in the GPT-4 ones. 

\begin{figure}[h]
\centering
\includegraphics[width=0.8\textwidth]{Plots/cbt.png}
\caption{Performance on the Children's Book Test dataset}
\label{fig:cbt}
\centering
\end{figure}

\subsection{FLOPs Analysis}

Observing the total training FLOPs is interesting as well. Ideally, each model would use the same number of FLOPs, since the same compute time on the same GPU was provided. However, from table \autoref{tab:flops}, this is not the case. Further investigation should that efficiency of MFU for each model is highly GPU dependent, and training on a different GPU using a similar compute budget may be more efficient on different models. The 8-layer model which uses more hidden layers in the FFN was less efficient than the baseline model, while being able to achieve very similar results. If we were to use FLOPs as the budget constraint, this model would have performed the best. 

\begin{table}[!h]
    \centering
    \begin{tabular}{rr}
        \hline
        Config & Total Training FLOPs\\
        \hline
        Baseline & $3.52 \times 10^{16}$ \\
        4-Layer  & $3.50 \times 10^{16}$ \\
        12-Layer & $3.53 \times 10^{16}$ \\
        Token Width Bias & $3.28 \times 10^{16}$ \\
        FFN Bias & $3.19 \times 10^{16}$ \\
        8-Layer Character & $3.41 \times 10^{16}$ \\
        5\% Character & $3.45 \times 10^{16}$ \\
        5\% BPE & $3.50 \times 10^{16}$ \\
        Mid-Training & $3.54 \times 10^{16}$\\
        \hline
    \end{tabular}
    \vspace{1em}
    \caption{Total FLOPs in 30min for each configuration}
    \label{tab:flops}
\end{table}

\section{Conclusion}

We show that given this 30 minutes constraint on a 3060ti, which is equivalent to roughly $3.5 \times 10^{16}$ training FLOPs, an 8-layer model with greater emphasize given to the FFN layer parameter count, using BPE tokenization gives the best performance. From the results when assigning the type of training data differently, we also show that the model has more biased towards the data it sees later on the training. 

\bibliographystyle{alpha}
\bibliography{references}

\appendix

\section{Other Figures}

\begin{figure}[h]
\centering
\includegraphics[width=0.8\textwidth]{Plots/vali_all.png}
\caption{Validation loss, recorded at intervals}
\label{fig:vali}
\centering
\end{figure}

\begin{figure}[h]
\centering
\includegraphics[width=0.8\textwidth]{Plots/test_all.png}
\caption{Final test loss}
\label{fig:test}
\centering
\end{figure}

\end{document}
